\documentclass[11pt,a4paper]{article}
\usepackage[utf8]{inputenc}
\usepackage[T1]{fontenc}
\usepackage[portuguese]{babel}
\usepackage{xcolor}
\usepackage{hyperref}
\usepackage{geometry}
\usepackage{listings}
\usepackage{tcolorbox}
\geometry{margin=2.5cm}
\definecolor{primary}{RGB}{41,128,185}
\definecolor{secondary}{RGB}{44,62,80}
\definecolor{codebg}{RGB}{240,240,240}
\lstset{backgroundcolor=\color{codebg},basicstyle=\ttfamily\small,breaklines=true,frame=single}
\title{\color{secondary}\Huge\textbf{Comunicação entre Pods e Services}}
\author{\color{primary}DevOps com Kubernetes -- 2026}
\date{}
\begin{document}
\maketitle


\section*{\color{primary}O que você aprenderá nesta página}
\begin{itemize}
\item Explicar como os Pods se comunicam dentro do cluster
\item Criar Services para estabilizar acesso a aplicações
\item Diferenciar ClusterIP, NodePort, LoadBalancer e ExternalName
\end{itemize}

Quando trabalhamos apenas com um \texttt{Deployment}, cada Pod criado recebe seu próprio endereço IP. Esse IP, porém, não deve ser usado como contrato de comunicação: Pods são substituídos, recriados e movidos entre nós. A solução do Kubernetes para esse problema é o \texttt{Service}, um recurso que oferece um ponto de acesso estável para um conjunto de Pods selecionados por \texttt{labels}.

O \texttt{Service} separa quem consome a aplicação de onde os Pods estão rodando naquele momento. Em vez de chamar diretamente o IP de um Pod, outro componente chama o nome DNS do Service, como \texttt{backend-svc.default.svc.cluster.local}. O Kubernetes mantém os \texttt{Endpoints} atualizados conforme os Pods entram ou saem do conjunto selecionado.

\section*{\color{primary}Labels, selectors e Endpoints}

Um Service encontra seus Pods por meio de seletores. Por isso, labels bem definidas deixam de ser detalhe estético e viram parte da arquitetura da aplicação. Se o seletor do Service não combinar com os labels do Pod, o Service existirá, mas não encaminhará tráfego para lugar nenhum.

\begin{lstlisting}[language=bash]
apiVersion: apps/v1
kind: Deployment
metadata:
  name: app-dep
spec:
  replicas: 3
  selector:
    matchLabels:
      app: app
  template:
    metadata:
      labels:
        app: app
    spec:
      containers:
        - name: app
          image: nginx:1.27
          ports:
            - containerPort: 80
---
apiVersion: v1
kind: Service
metadata:
  name: app-svc
spec:
  selector:
    app: app
  ports:
    - port: 80
      targetPort: 80
  type: ClusterIP
\end{lstlisting}

Depois de aplicar o manifesto, os comandos \texttt{kubectl get svc}, \texttt{kubectl get endpoints app-svc} e \texttt{kubectl describe svc app-svc} ajudam a verificar se o Service realmente encontrou os Pods.

\section*{\color{primary}Tipos de Service}

O tipo \texttt{ClusterIP} é o padrão e expõe a aplicação apenas dentro do cluster. Ele é ideal para comunicação interna entre microsserviços. O \texttt{NodePort} abre uma porta em todos os nós do cluster e encaminha o tráfego para o Service. O \texttt{LoadBalancer} solicita ao provedor de nuvem um balanceador externo. O \texttt{ExternalName} cria um alias DNS para um serviço externo.

\section*{\color{primary}Headless Service}

Um Headless Service é criado com \texttt{clusterIP: None}. Nesse caso, o Kubernetes não fornece um IP virtual único para balanceamento. Em vez disso, o DNS retorna os endereços dos Pods diretamente. Esse padrão é importante para aplicações que precisam conhecer a identidade individual de cada réplica, como bancos de dados e sistemas distribuídos.

\begin{lstlisting}[language=bash]
apiVersion: v1
kind: Service
metadata:
  name: postgres-headless
spec:
  clusterIP: None
  selector:
    app: postgres
  ports:
    - port: 5432
\end{lstlisting}

\section*{\color{primary}Boas práticas}

Use nomes claros, mantenha labels consistentes e evite depender de IPs de Pods. Para comunicação interna, comece com \texttt{ClusterIP}. Para publicação externa, prefira Ingress, Gateway API ou um LoadBalancer controlado. Sempre valide se os Endpoints aparecem.

\end{document}
